% =========================================================================
\section{Distribuciones de dipolos: la polarización}\label{sec:polarizacion}
% =========================================================================
Del mismo modo que una distribución continua de carga es el paso al continuo de un sistema de cargas puntuales, $\rho=\dif q/\dVp$,
podemos considerar una distribución continua de \emph{dipolos}. Es la situación de un material dieléctrico, en el que un número enorme de
dipolos moleculares, muy próximos entre sí, ocupa una región $V'$ del espacio.

% -------------------------------------------------------------------------
\subsection{Densidad de momento dipolar}\label{sec:pol-def}
% -------------------------------------------------------------------------
\begin{definicion}[Polarización]\label{def:polarizacion}
Sea una región $V'$ con una distribución continua de dipolos. Se llama \emph{polarización} al campo vectorial
\begin{equation}\label{eq:polarizacion}
  \vP(\rp)=\frac{\dif\vp}{\dVp},\qquad\text{es decir,}\qquad \dif\vp=\vP(\rp)\,\dVp ,
\end{equation}
momento dipolar por unidad de volumen en el punto fuente $\rp$. Su unidad es $[\vP]=\mathrm{C\,m}/\mathrm{m^{3}}=\mathrm{C/m^{2}}$, la misma que la de una
densidad superficial de carga.
\end{definicion}

\begin{observacion}[Analogía con la densidad de carga]\label{obs:analogia}
\begin{center}\vspace{2pt}
\renewcommand{\arraystretch}{1.25}
\begin{tabular}{@{}lcc@{}}
  \toprule
  & \textsc{Carga puntual} & \textsc{Dipolo puntual}\\
  \midrule
  Magnitud discreta & $q$ & $\vp$\\
  Densidad continua & $\rho(\rp)=\dif q/\dVp$ & $\vP(\rp)=\dif\vp/\dVp$\\
  Unidad de la densidad & $\mathrm{C/m^{3}}$ & $\mathrm{C/m^{2}}$\\
  Potencial elemental & $\dfrac{\dif q}{4\pi\eo R}$ & $\dfrac{\dif\vp\cdot\uR}{4\pi\eo R^{2}}$\\
  \bottomrule
\end{tabular}
\end{center}
\vspace{2pt}
La polarización es una magnitud fundamental en el estudio de los materiales dieléctricos, que abordaremos más adelante en el curso.
\end{observacion}

% =========================================================================
% Figura (apuntes manuscritos, pág. 5, y de teoría, pág. II-4): región V'
% polarizada; el elemento dV' en r' contiene el dipolo dp = P dV', que crea
% en P (posición r) el potencial dphi = dp.u_R/(4 pi eps0 R^2).
% =========================================================================
\begin{figure}[ht]
\centering
\begin{tikzpicture}[x=1cm, y=1cm, >=Stealth]
  \ejesO{4.4}{6.6}{1.6}
  \regionV
  % pequeños dipolos (polarización) en el interior
  \foreach \px/\py in {1.45/2.5, 1.75/3.35, 2.55/3.55, 3.15/3.1, 3.2/2.25, 2.45/1.95, 1.75/1.95}{%
    \draw[-{Stealth[length=2.6pt,width=2.2pt]}, magacero!75, line width=0.55pt] (\px-0.12,\py-0.08) -- (\px+0.14,\py+0.1);}
  \node[etiq, text=magprof] at (0.8,3.65) {$V'$};
  % elemento de volumen
  \coordinate (Q) at (2.35,2.75);
  \coordinate (P) at (6.2,3.25);
  \draw[vect] (0,0) -- ($(Q)+(-0.1,-0.12)$) node[pos=0.5, above left=-1pt, etiq] {$\rp$};
  \cubito{Q}
  \node[etiq, font=\footnotesize, anchor=south east] at ($(Q)+(-0.1,0.12)$) {$\dVp$};
  \draw[vect] (0,0) -- (P) node[pos=0.72, below right=-2pt, etiq] {$\vc{r}$};
  \draw[vectrojo] ($(Q)+(0.22,0.06)$) -- (P) node[pos=0.6, above=2pt, etiq, text=rojoteo] {$\vc{R}=\vc{r}-\rp$};
  \fill[magprof] (P) circle (1.5pt) node[right=2pt, etiq] {$P$};
  % normal exterior
  \coordinate (B) at (3.3,1.8);
  \draw[vectverde] (B) -- ++(0.42,-0.44) node[below right=-2pt, etiq, text=verdeteo] {$\un$};
  \node[etiq, font=\footnotesize, text=magacero, anchor=south] at (2.45,3.98) {$S'=\partial V'$};
  \fill[black] (0,0) circle (1.2pt) node[below right=1pt, etiq, font=\footnotesize] {$O$};
  % recuadro con las relaciones
  \node[cuadro, anchor=north west, align=left, font=\footnotesize, inner sep=4pt] at (7.0,2.6)
    {$\dif\vp=\vP(\rp)\,\dVp$\\[4pt]
     $\dif\phi(\vc{r})=\dfrac{\dif\vp\cdot\uR}{4\pi\eo R^{2}}$\\[7pt]
     $\rhoP=-\nabp\cdot\vP$,\quad $\sigP=\vP\cdot\un$};
\end{tikzpicture}
\caption{Distribución continua de dipolos en una región $V'$ limitada por la superficie $S'$, de normal exterior $\un$. El elemento $\dVp$ situado en $\rp$
contiene el momento dipolar $\dif\vp=\vP\,\dVp$, que contribuye al potencial en el punto campo $P$ como un dipolo puntual a la distancia $R$. El potencial
total equivale al de las cargas de polarización $\rhoP$ y $\sigP$ (\cref{teo:cargas-pol}).}
\label{fig:polarizacion}
\end{figure}


% -------------------------------------------------------------------------
\subsection{Potencial de una distribución de dipolos}\label{sec:pol-potencial}
% -------------------------------------------------------------------------
Cada elemento de volumen $\dVp$ situado en $\rp$ contiene un dipolo $\dif\vp=\vP\,\dVp$ que, por la \cref{prop:dipolo-Rp}, crea en $\vc{r}$ el potencial
$\dif\phi=\dif\vp\cdot\uR/(4\pi\eo R^{2})$ (\cref{fig:polarizacion}). Por superposición,
\begin{equation}\label{eq:pot-polarizacion}
  \phi(\vc{r})=\frac{1}{4\pi\eo}\int_{V'}\frac{\vP(\rp)\cdot\uR}{R^{2}}\,\dVp=\frac{1}{4\pi\eo}\int_{V'}\frac{\vP(\rp)\cdot\vc{R}}{R^{3}}\,\dVp .
\end{equation}
Esta integral tiene una interpretación sorprendente: el potencial es el mismo que crearía una cierta distribución de \emph{carga}, en volumen y en la
superficie que limita $V'$.

\begin{lema}[Gradiente de $1/R$ respecto de las coordenadas fuente]\label{lem:grad-1R}
Para $\vc{r}\neq\rp$,
\begin{equation}\label{eq:grad-1R}
  \nabp\Bigl(\frac{1}{R}\Bigr)=\frac{\vc{R}}{R^{3}}=\frac{\uR}{R^{2}}=-\nab\Bigl(\frac{1}{R}\Bigr).
\end{equation}
\end{lema}
\begin{proof}
Es el gradiente de $R^{n}$ con $n=-1$ estudiado en el Tema~0~\cite{TEMA0}: $\nab R^{n}=nR^{n-1}\uR$, de donde $\nab(1/R)=-\uR/R^{2}$. Como
$\vc{R}=\vc{r}-\rp$, derivar respecto de $x'_j$ equivale a derivar respecto de $x_j$ cambiando el signo, luego $\nabp(1/R)=+\uR/R^{2}$.
\end{proof}

\begin{teorema}[Cargas de polarización]\label{teo:cargas-pol}
Sea $\vP$ una polarización de clase $C^{1}$ en una región $V'$ limitada por una superficie cerrada $S'=\partial V'$ con normal unitaria exterior
$\un$. El potencial \cref{eq:pot-polarizacion} es
\begin{equation}\label{eq:pot-cargas-pol}
  \phi(\vc{r})=\frac{1}{4\pi\eo}\int_{V'}\frac{\rhoP(\rp)}{R}\,\dVp+\frac{1}{4\pi\eo}\oiint_{S'}\frac{\sigP(\rp)}{R}\,\dSp,
\end{equation}
donde
\begin{equation}\label{eq:densidades-pol}
  \boxed{\;\rhoP=-\nabp\cdot\vP,\qquad \sigP=\vP\cdot\un\;}
\end{equation}
son las \emph{densidades volúmica y superficial de carga de polarización}.
\end{teorema}
\begin{proof}
Por el \cref{lem:grad-1R}, el integrando de \cref{eq:pot-polarizacion} es $\vP\cdot\nabp(1/R)$. La regla del producto para la divergencia,
$\nabp\cdot(f\vc{A})=f\,\nabp\cdot\vc{A}+\vc{A}\cdot\nabp f$, con $f=1/R$ y $\vc{A}=\vP$, da
\[
  \nabp\cdot\Bigl(\frac{\vP}{R}\Bigr)=\frac{\nabp\cdot\vP}{R}+\vP\cdot\nabp\Bigl(\frac{1}{R}\Bigr)
  \quad\Longrightarrow\quad
  \frac{\vP\cdot\uR}{R^{2}}=-\frac{\nabp\cdot\vP}{R}+\nabp\cdot\Bigl(\frac{\vP}{R}\Bigr).
\]
Integrando en $V'$,
\[
  \phi(\vc{r})=\frac{1}{4\pi\eo}\Biggl[\int_{V'}\frac{-\nabp\cdot\vP}{R}\,\dVp+\int_{V'}\nabp\cdot\Bigl(\frac{\vP}{R}\Bigr)\dVp\Biggr],
\]
y por el teorema de la divergencia (de Gauss)~\cite{TEMA0}, con $\dif\vc{S}\,'=\dSp\,\un$, la segunda integral es
$\oiint_{S'}\frac{\vP\cdot\un}{R}\,\dSp$. Basta llamar $\rhoP$ y $\sigP$ a los numeradores. (Si $\vc{r}\in V'$, la singularidad en $R=0$ es integrable y
el resultado se mantiene excluyendo una pequeña esfera alrededor de $\vc{r}$ y haciendo tender su radio a cero.)
\end{proof}

\begin{observacion}[Cargas ligadas]\label{obs:ligadas}
Comparando \cref{eq:pot-cargas-pol} con el potencial de una distribución de carga (\cref{eq:integrales}), una distribución de dipolos equivale, a efectos de
potencial y de campo, a una carga $\rhoP$ en volumen más una carga $\sigP$ en superficie. Estas cargas no son libres: proceden del desplazamiento
microscópico de las cargas ligadas de cada molécula, y por eso la terminología habitual de los dieléctricos las llama \emph{cargas ligadas}. La normal
$\un$ apunta, por convenio, hacia fuera del volumen. Físicamente, $\sigP$ es la carga que «asoma» en la superficie al no estar compensada por la
molécula vecina, y $\rhoP$ aparece donde la polarización no es uniforme (si $\vP$ es constante, $\rhoP=0$ y solo hay carga superficial).
\end{observacion}

\begin{proposicion}[La carga total de polarización es nula]\label{prop:carga-pol-nula}
\begin{equation}\label{eq:carga-pol-nula}
  Q_{P}=\int_{V'}\rhoP\,\dVp+\oiint_{S'}\sigP\,\dSp=0 .
\end{equation}
\end{proposicion}
\begin{proof}
Por el teorema de la divergencia,
\[
  Q_{P}=-\int_{V'}\nabp\cdot\vP\,\dVp+\oiint_{S'}\vP\cdot\un\,\dSp=-\oiint_{S'}\vP\cdot\dif\vc{S}\,'+\oiint_{S'}\vP\cdot\dif\vc{S}\,'=0 .
\]
Es lo esperable: cada dipolo es neutro, y por tanto también lo es cualquier conjunto de ellos.
\end{proof}
